% ======================================================================
% CHAPTER 6: NETWORK TESTBED TOPOLOGIES & DYNAMIC WORKLOAD SIMULATION
% ======================================================================
\chapter{Network Testbed Topologies \& Dynamic Workload Simulation}
\label{chap:topologies}

\section{Evaluated Network Testbed Topologies}
To evaluate the adaptability, scalability, and generalizability of the developed reinforcement learning traffic engineering platform, the system implements and evaluates five topologically diverse network architectures via \texttt{topology/topology\_library.py}. These topologies span multi-tenant cloud data center fabrics, enterprise campus backbones, and continental-scale Wide Area Networks (WANs).

\subsection{Hierarchical Tree Fabric (7 Switches, 8 Hosts)}
The Hierarchical Tree topology represents the classical 3-tier enterprise campus design:
\begin{itemize}
    \item \textbf{Core Tier}: Switch $s_1$ acts as the single central transit router.
    \item \textbf{Aggregation Tier}: Switches $s_2$ and $s_3$ aggregate traffic from edge switches.
    \item \textbf{Edge Tier}: Switches $s_4, s_5, s_6, s_7$ connect directly to end hosts $h_1$ through $h_8$.
    \item \textbf{Lateral Cross-Links}: Redundant horizontal bypass links connect adjacent edge switches ($s_4 \leftrightarrow s_6$ and $s_5 \leftrightarrow s_7$) with $100\text{ Mbps}$ capacity.
\end{itemize}
Under traditional shortest-path routing, all inter-pod traffic is funneled through core switch $s_1$, creating severe core bottlenecks while the lateral cross-links remain completely unutilized. The Hierarchical Tree provides a clean testbed for verifying lateral link offloading.

\subsection{Fat-Tree Clos Fabric ($k=4$, 20 Switches, 16 Hosts)}
The Fat-Tree topology \cite{al2008scalable} implements a multi-stage Clos network widely deployed in modern private and public cloud data centers:
\begin{itemize}
    \item \textbf{Pod Architecture}: Four distinct pods ($k=4$). Each pod contains two aggregation switches and two edge switches.
    \item \textbf{Core Tier}: Four centralized core switches interconnected with all aggregation switches:
    \begin{equation}
        N_{\text{core}} = (k/2)^2 = (4/2)^2 = 4 \text{ Core Switches}
    \end{equation}
    \item \textbf{Edge and Host Connectivity}: Eight edge switches connect to sixteen bare-metal hosts ($h_1$ through $h_{16}$).
    \item \textbf{Bisection Bandwidth}: Possesses full bisection bandwidth: any arbitrary bisection cut intersects 16 full-rate links.
\end{itemize}
Because multiple parallel paths exist between any two pods, the Fat-Tree fabric presents the ideal environment to evaluate multi-path load balancing and elephant flow collision avoidance.

\subsection{Abilene US Backbone Network (12 Switches, 8 Hosts)}
The Abilene network represents the high-speed continental backbone of the Internet2 nationwide research network across the continental United States:
\begin{itemize}
    \item \textbf{Topological Structure}: 12 switches interconnected by 30 directed transmission links spanning Seattle, Sunnyvale, Los Angeles, Denver, Kansas City, Houston, Chicago, Indianapolis, Atlanta, Washington DC, and New York.
    \item \textbf{Asymmetric Propagation Delays}: Unlike local data centers (where delay is negligible), Abilene links exhibit real-world geographic propagation delays ranging from $2.0\text{ ms}$ (Chicago to Indianapolis) to $8.5\text{ ms}$ (Seattle to Sunnyvale).
\end{itemize}
Abilene provides the critical benchmark for evaluating the trade-off between detour path length (WAN propagation latency) and core link queuing delay.

\subsection{NSFNet Continental Mesh (14 Switches, 8 Hosts)}
The National Science Foundation Network (NSFNet) was the historic foundation of the modern Internet. It represents an irregular, non-symmetric continental mesh:
\begin{itemize}
    \item \textbf{Topological Structure}: 14 switches interconnected by 42 directed transmission links.
    \item \textbf{Node Degree Diversity}: Switch node degrees range from 2 (e.g., edge transit nodes) up to 4 (e.g., highly connected central transit hubs).
\end{itemize}
The irregular connectivity of NSFNet tests the RL policy's capacity to navigate complex meshed bottlenecks without uniform geometric symmetry.

\subsection{Spine-Leaf Data Center Fabric (12 Switches, 16 Hosts)}
The Spine-Leaf architecture represents the modern state-of-the-art for hyperscale cloud data center fabrics (e.g., AWS, Azure, Google Cloud):
\begin{itemize}
    \item \textbf{Spine Tier}: 4 independent high-capacity spine switches ($s_1, s_2, s_3, s_4$).
    \item \textbf{Leaf Tier}: 8 access leaf switches ($s_5$ through $s_{12}$), each directly attached to compute servers.
    \item \textbf{Bipartite Interconnection}: Every leaf switch is connected to every single spine switch, creating a complete bipartite graph ($K_{4, 8}$).
\end{itemize}
Under classical single-path routing, a leaf pair picks a single spine, leaving three-quarters of the fabric bandwidth wasted. The Spine-Leaf fabric rigorously evaluates the agent's ability to achieve 100\% fabric dispersion.

\section{Topological Graph Metrics Comparative Analysis}
Table~\ref{tab:graph_metrics_comparative} presents the rigorous graph-theoretic metrics for all five benchmark fabrics.

\begin{table}[htbp]
\centering
\caption{Comprehensive Graph-Theoretic Metrics of Evaluated Network Topologies}
\label{tab:graph_metrics_comparative}
\small
\begin{tabularx}{\textwidth}{l c c c c c c c}
\toprule
\textbf{Topology Name} & \textbf{$|V|$} & \textbf{$|E|$} & \textbf{Hosts} & \textbf{Diameter} & \textbf{Avg Degree} & \textbf{Bisection BW} & \textbf{Clustering Coeff} \\
\midrule
\textbf{Hierarchical Tree} & 7 & 16 & 8 & 4 hops & 2.29 & 200 Mbps & 0.285 \\
\textbf{Fat-Tree ($k=4$)} & 20 & 64 & 16 & 6 hops & 3.20 & 1,600 Mbps & 0.000 \\
\textbf{Abilene US Backbone} & 12 & 30 & 8 & 5 hops & 2.50 & 400 Mbps & 0.121 \\
\textbf{NSFNet Continental} & 14 & 42 & 8 & 4 hops & 3.00 & 600 Mbps & 0.244 \\
\textbf{Spine-Leaf Fabric} & 12 & 64 & 16 & 2 hops & 5.33 & 3,200 Mbps & 0.000 \\
\bottomrule
\end{tabularx}
\end{table}

\section{Synthetic Traffic Generation Engine}
Realistic workload modeling is essential to prevent reinforcement learning agents from overfitting to trivial static traffic patterns. The synthetic traffic generation engine (\texttt{simulate\_live\_traffic.py} and \texttt{topology/traffic\_generator.sh}) implements a multi-tier stochastic traffic injection framework.

\subsection{Poisson Process for Background Mice Flows}
Background transactional traffic (DNS queries, HTTP micro-requests, RPC handshakes) is modeled as an independent Poisson arrival process. The probability of $k$ flow arrivals occurring within time interval $t$ is:
\begin{equation}
    \mathbb{P}(N(t) = k) = \frac{(\lambda_{\text{mice}} \cdot t)^k \cdot e^{-\lambda_{\text{mice}} t}}{k!}
\end{equation}
where $\lambda_{\text{mice}} = 15.0\text{ flows/second}$. Each mice flow is allocated a modest bandwidth demand uniformly sampled:
\begin{equation}
    D_{\text{mice}} \sim \mathcal{U}(3.0\text{ Mbps}, \, 8.0\text{ Mbps})
\end{equation}
Mice flows persist for short durations $T_{\text{mice}} \sim \mathcal{U}(1.5\text{s}, \, 4.0\text{s})$ before terminating and releasing link buffer capacity.

\subsection{Heavy-Tailed Pareto Distribution for Elephant Microbursts}
Bulk data transfers, database backups, and VM live migrations exhibit heavy-tailed size distributions, where a tiny fraction of flows account for the vast majority of transmitted octets. We model elephant flows via the classical \textbf{Bounded Pareto Distribution}:
\begin{equation}
    f(x; \alpha, x_m, x_M) = \frac{\alpha \cdot x_m^\alpha \cdot x^{-\alpha - 1}}{1 - (x_m / x_M)^\alpha} \quad \text{for } x_m \le x \le x_M
\end{equation}
where:
\begin{itemize}
    \item Shape parameter $\alpha = 1.2$ reflects heavy-tailed Internet traffic measurements.
    \item Minimum flow demand $x_m = 35.0\text{ Mbps}$.
    \item Maximum flow demand $x_M = 85.0\text{ Mbps}$.
\end{itemize}
Elephant flows are injected between diagonal edge switches across pod boundaries, persisting for $T_{\text{elephant}} \sim \mathcal{U}(10.0\text{s}, \, 25.0\text{s})$, designed to aggressively saturate standard shortest paths.

\subsection{Spatial Hotspot Clustered Surges}
In enterprise campus and multi-tenant networks, traffic frequently exhibits spatial locality (e.g., sudden video broadcast from a conference hall, distributed MapReduce shuffle phase). We model this via spatial hotspot clustering:
\begin{equation}
    \mathbb{P}(\text{Source} = v_i, \, \text{Dest} = v_j) = \begin{cases}
        0.75 \cdot \frac{1}{|V_{\text{hotspot}}|^2} & \text{if } v_i, v_j \in V_{\text{hotspot}} \\[0.5em]
        0.25 \cdot \frac{1}{|V \setminus V_{\text{hotspot}}|^2} & \text{otherwise}
    \end{cases}
\end{equation}
where $V_{\text{hotspot}}$ represents a localized cluster of edge switches (e.g., switches $s_4$ and $s_5$).

\subsection{Targeted Core Jamming Stress Model}
To stress-test emergency anti-congestion rerouting, the engine implements targeted core link pre-loading. Prior to injecting measurement flows, central transit links are pre-loaded to $85\%$--$98\%$ utilization by generating persistent synthetic saturation streams:
\begin{equation}
    \text{PreLoad}(e_{\text{core}}) \sim \mathcal{U}(0.85 \cdot C_e, \, 0.98 \cdot C_e)
\end{equation}
This forces the RL agent to immediately detect the asymptotic barrier penalty $\Phi_{\text{cong}}$ and dynamically steer new arrivals onto lateral bypass paths.

\section{The 3-Phase Curriculum Training Strategy}
Training an RL agent directly on high-intensity multi-commodity traffic causes policy instability: the agent encounters severe penalties before understanding basic path navigation. To ensure smooth, monotonic policy convergence, we designed a \textbf{3-Phase Curriculum Training Loop} executed over 1,000 episodes (\texttt{benchmark\_evaluation.py}):

\begin{figure}[htbp]
\centering
\begin{minipage}{0.92\textwidth}
\centering
\fbox{\parbox{0.95\textwidth}{
\vspace{0.2cm}
\centering
{\large \textbf{3-Phase Curriculum Training Architecture}}\\[0.3cm]
\begin{tabular}{c c c}
\fbox{\parbox{0.28\textwidth}{\centering \textbf{Phase 1: Exploration}\\ \footnotesize Episodes 1--250\\ $\epsilon: 1.0 \to 0.5$\\ Light Poisson Mice Flows\\ Discovery of Path Delays}} 
& $\xrightarrow{\text{Stabilization}}$ 
& \fbox{\parbox{0.28\textwidth}{\centering \textbf{Phase 2: Learning \& Stress}\\ \footnotesize Episodes 251--750\\ $\epsilon: 0.5 \to 0.01$\\ Bursty Pareto Elephants\\ Core Jamming Injections}} \\
& & $\Downarrow$ \footnotesize Convergence \\
& & \fbox{\parbox{0.28\textwidth}{\centering \textbf{Phase 3: Exploitation}\\ \footnotesize Episodes 751--1,000\\ $\epsilon = 0.01$ (Floor)\\ Steady-State Evaluation\\ Policy Fine-Tuning}} \\
\end{tabular}
\vspace{0.2cm}
}}
\end{minipage}
\caption{3-Phase Curriculum Training Progression over 1,000 Episodes.}
\label{fig:curriculum_progression}
\end{figure}

\subsection{Phase 1: Pure Exploration Baseline (Episodes 1--250)}
\begin{itemize}
    \item \textbf{Exploration Rate}: $\epsilon$ decays linearly from $1.0$ down to $0.50$.
    \item \textbf{Workload}: Low-intensity Poisson background mice flows (15--25\% average link load).
    \item \textbf{Objective}: The agent explores the action space uniformly, discovering candidate path hop counts and baseline physical propagation delays. The PER buffer is populated with diverse baseline transitions.
\end{itemize}

\subsection{Phase 2: Stress \& Bottleneck Rerouting (Episodes 251--750)}
\begin{itemize}
    \item \textbf{Exploration Rate}: $\epsilon$ anneals from $0.50$ down to the minimum floor $\epsilon_{\text{min}} = 0.01$.
    \item \textbf{Workload}: Heavy Pareto elephant flows, asymmetric pod surges, and core jamming injections (70--98\% link load).
    \item \textbf{Objective}: The agent repeatedly encounters the steep non-linear barrier penalty $\Phi_{\text{cong}}$ on core links. Through TD-error backpropagation, the Q-network learns to associate saturated core paths with large negative returns, developing the proactive policy to steer traffic onto lateral bypasses before queue overflow occurs.
\end{itemize}

\subsection{Phase 3: Fine-Tuning \& Steady Exploitation (Episodes 751--1000)}
\begin{itemize}
    \item \textbf{Exploration Rate}: Fixed at the minimum exploration floor $\epsilon = 0.01$ ($99\%$ greedy exploitation).
    \item \textbf{Workload}: Mixed heterogeneous production traffic profiles rotating across all five network topologies.
    \item \textbf{Objective}: The policy stabilizes at the optimal multi-objective Pareto boundary. Smooth L1 loss plateaus at $0.45$--$1.36$, and episodic return achieves positive convergence ($+4.16$).
\end{itemize}
